\documentclass[12pt,a4paper]{moderncv}
\moderncvtheme[blue]{classic}
\usepackage[utf8]{inputenc}
\usepackage[scale=0.76]{geometry}

% Footer
\usepackage{fancyhdr}
\pagestyle{fancy}
\lfoot{Cristián Maureira-Fredes | \cc{2018}}

% Commands
\newcommand{\myname}{\textbf{Cristián Maureira}}
\newcommand{\myeman}{\textbf{Maureira, Cristián}}
\newcommand{\mynamefull}{\textbf{C. Maureira-Fredes}}
\newcommand{\utfsm}{Universidad Técnica Federico Santa María (UTFSM)}
\newcommand{\valpo}{Valparaíso, Chile}
\newcommand{\gray}{\textcolor{gray}}
\newcommand{\myurl}[1]{\gray{\url{#1}}}
\newcommand{\aei}{Max-Planck-Institut für Gravitationsphysik (Albert-Einstein-Institut)}
\newcommand{\cc}{\textcolor{color1}}

\definecolor{mred}{RGB}{225,67,46}
\definecolor{mblue}{RGB}{75,92,105}
\definecolor{mgreen}{RGB}{121,189,154}
\definecolor{mcyan}{RGB}{70,197,218}
\definecolor{mwhite}{RGB}{221,220,215}
\definecolor{mgray}{RGB}{143,143,143}
\definecolor{mpetrol}{RGB}{59,134,134}
\definecolor{mnight}{RGB}{0,180,204}

\newcommand{\mred}{\textcolor{mred}}
\newcommand{\mblue}{\textcolor{mblue}}
\newcommand{\mgreen}{\textcolor{mgreen}}
\newcommand{\mcyan}{\textcolor{mcyan}}
\newcommand{\mwhite}{\textcolor{mwhite}}
\newcommand{\mgray}{\textcolor{mgray}}
\newcommand{\mpetrol}{\textcolor{mpetrol}}
\newcommand{\mnight}{\textcolor{mnight}}

% Font
\usepackage{fontspec}
\setmainfont{Fira Sans Light}
%\setmainfont{Myriad Pro}

% Left column lenght
\setlength{\hintscolumnwidth}{2.2cm}

%%%%%%%%%%%%%%%%%%%%%%%%%%%%%%%%%%%%%%%%%%%% Personal Information
\firstname{\huge Cristián}
\familyname{\huge Maureira-Fredes}
\title{\Large M.Sc. Computer Science,\newline
Informatics Engineer}
%\address{Gerdeschützenweg 11B, D-12203}{Berlin, Germany}
%\mobile{+49 152 5695 6241}
%\phone{+49 331 567 7323}
%\fax{fax (optional)}
\email{cmaureirafredes@gmail.com}
\homepage{maureira.xyz}
\extrainfo{
\includegraphics[width=0.023\textwidth]{icon/Linkedin.pdf}\
\href{https://www.linkedin.com/in/cmaureir}{Cristián Maureira Fredes}\\
\includegraphics[width=0.023\textwidth]{icon/GitLab_Logo.pdf}\
\href{http://gitlab.com/u/cmaureir}{https://gitlab.com/cmaureir}\\
\includegraphics[width=0.023\textwidth]{icon/GitHub_Logo.pdf}\
\href{http://github.com/cmaureir}{https://github.com/cmaureir}}
%\photo[82pt]{photo/me}

%%%%%%%%%%%%%%%%%%%%%%%%%%%%%%%%%%%%%%%%%%%% Document
\begin{document}
\maketitle

%(Earliest possible entry 01.01.2018)
%%%%%%%%%%%%%%%%%%%%%%%%%%%%%%%%%%%%%%%%%%%% Education
\section{\LARGE Education}

\cventry{\cc{\textbf{2015 to\\ present}}}
        {Doctoral Student}
        {Fakultät für Mathematik und Physik}
        {Leibniz Universität Hannover}
        {Hannover}
        {Germany}
\cventry{\cc{\textbf{2013 to\\ present}}}
        {Doctoral Student}
        {}
        {\aei}
        {Potsdam-Golm}
        {Germany}
\cventry{2011 to\\2013}
        {M.Sc. Computer Science}
        {}
        {{\utfsm}}
        {Valparaíso}
        {Chile}
\cventry{2006 to\\ 2013}
        {Informatics Engineering}
        {}
        {{\utfsm}}
        {Valparaíso}
        {Chile}
%\cventry{1993-2005}
%        {Elementary and High School}
%        {Instituto del Puerto}
%        {San Antonio, Valparaíso}
%        {}{}

%\vspace{-0.5cm}

%%%%%%%%%%%%%%%%%%%%%%%%%%%%%%%%%%%%%%%%%%%% Experience
\section{\LARGE Experience}
\subsection{\Large Vocational}
\vspace{0.3cm}

\cventry{Current}
        {Doctoral studies}
        {Working in the numerical modelling of \cc{Dense Stellar systems}
        and \cc{Super Massive Black Hole} binaries.\newline
        My research consist mainly in astrophysical simulations using \cc{HPC}
        techniques and facilities, to analyse the evolution and properties
        of different systems.\newline
        Most of the \cc{analysis} is done processing large simulation data sets,
        from which a proper extraction of useful astrophysical and statistical
        information is crucial.
        }
        {\aei}
        {Germany}
        {\myurl{http://aei.mpg.de}}

\newpage
\cventry{May 2012 to\\ Nov 2012}
        {Master Student (Visitor)}
        {Thesis development, %
        ``A modular, GPU-based, direct-summation $N-$body integrator'' }
        {\aei}
        {Germany}
        {\myurl{http://aei.mpg.de}}

\cventry{2011 to\\ 2012}
        {CTI-HPC Research Assistant}
        {Researching %
         Astrophysical simulations using GPU technologies, %
         High Frequency Trading using CUDA and MKL in Neural Networks evaluation %
         and Image processing and analysis, among others HPC-related projects.\newline
         Officially named a CUDA Teaching Center by NVIDIA.}
        {Center for Technological Innovation in High Performance Computing (CTI-HPC)}
        {UTFSM, Chile}
        {}

\cventry{2011 to\\ 2012}
        {ALMA-UTFSM Representative and developer}
        {In charge of the national and international representation of the group,
         coordinating projects in collaboration with the ALMA observatory
         (including ESO and the NRAO) at the UTFSM.}
        {Computer Systems Research Group (CSRG)}
        {UTFSM, Chile}
        {\myurl{http://www.csrg.cl/portal/alma-utfsm/}}

\cventry{2010 to\\ 2011}
        {ARM Research Group Member}
        {Research ARM architecture systems for the development of micro computers
         with low energy consumption}
        {Computer Systems Research Group (CSRG)}
        {UTFSM, Chile}
        {\myurl{http://www.csrg.cl/portal/arm/}}

\cventry{2010 to\\ 2011}
        {Assistantship}
        {Courses: Scientific Computing I and II}
        {Informatics Department}
        {UTFSM, Chile}
        {}

\cventry{Jul 2009 to\\ Dec 2012}
        {ALMA-UTFSM Research (Project leader)}
        {Project Leader of ``Software development, Analysis and Testing (SWAT)''
         Team, in charge of the development of the software at the ALMA-UTFSM group.
         Maintenance and development of the ``Sampling System GUI'' project,
         build upon the ALMA Common Software framework (ACS).
         Member of ``Alarms Configuration GUI'', with the propose of develop a
         Graphic User Interface to configure the ACS Alarms System.}
        {Computer Systems Research Group (CSRG)}
        {UTFSM, Chile}
        {\myurl{http://www.csrg.cl/portal/alma-utfsm/}}

\cventry{2008 to\\ 2009}
        {Project Leader}
        {$\mu$Bot Interface, A Generic Graphic User Interface to robotic platform
        control, a ``Vortex Labs'' project, PIE>A Sponsored Project.
        Application written in Python/C++ using Qt as graphical library.
        The software included a procedure editor, and a embedded python console
        to test the routines, XML profiles of each robotic system, etc.}
        {Computer Systems Research Group (CSRG)}
        {UTFSM, Chile}
        {\myurl{http://www.piea.usm.cl}}

\cventry{Aug 2008 to\\ Jul 2009}
        {ALMA-UTFSM Research (Assistantship)}
        {Member of ``Software development, Analysis and Testing'' (SWAT) Team,
        in charge of the development of the software at the ALMA-UTFSM group.
        Maintenance and development of the ``Sampling System GUI'' project,
        build upon the ALMA Common Software framework (ACS).
        Member of the ``Repackaging ACS for Embedded Systems project'' (RAES),
        with the purpose of separate the run-time and development
        proper of this build environment}
        {Computer Systems Research Group (CSRG)}
        {UTFSM, Chile}
        {\myurl{http://www.csrg.cl/portal/alma-utfsm/}}

\cventry{2008 to\\ 2009}
        {Project Member}
        {Rescue Live CD, A Fedora official Spin release,
         PIE>A, Sponsored Project}
        {Computer Systems Research Group (CSRG)}
        {UTFSM, Chile}
        {\myurl{http://www.piea.usm.cl}}

\cventry{Mar 2008 to\\ Dec 2008}
        {Assistantship}
        {Courses: Informatics Fundaments  I and II}
        {Informatics Department}
        {UTFSM, Chile}
        {}

\cventry{2007 to\\ 2008}
        {Computer Laboratory Assistantship}
        {Member of the ``Computer Laboratory'' (LabComp),
         laboratory in charge of all the Informatics Student services,
         performing the tasks of system administrator (networking, scripting and
         teaching Linux courses)}
        {UTFSM, Chile}
        {\myurl{http://labcomp.inf.utfsm.cl}}{}

\cventry{Jul 2007 to\\ Dec 2007}
        {Programming Laboratory Assistantship}
        {C language}
        {Informatics Department}
        {UTFSM, Chile}
        {}

\newpage
%%%%%%%%%%%%%%%%%%%%%%%%%%%%%%%%%%%%%%%%%%%% Internships
\section{\LARGE Experience}
\subsection{\Large Internships}
\vspace{0.3cm}

\cventry{2011}{Professional Internship}
        {{\utfsm}}
        {Center for Technological Innovation in High Performance Computing (CTI-HPC)}
        {Working on the implementation of a gravitational $N$-body simulations,
         using different parallel computing approaches, as POSIX threads, OpenMP
         and GPU Computing.}
        {UTFSM, Chile}
        {}

\cventry{2010}{Professional Internship}
        {{\utfsm}}
        {Working on the implementation of an alternative to the ALMA Common
         Software (ACS) Notification Channel using OpenDDS, \myurl{http://www.opendds.org/}}
        {UTFSM, Chile}
        {\myurl{http://csrg.cl}}

\cventry{2009}{Professional Internship}
        {Associated Universities Incorporated (AUI), Joint ALMA Observatory}
        {Working on the Java implementation of the ACS Basic Control Interface
         Specification; Developing a GUI tool to perform test in the Antenna Bus
         Master (ABM) and Teaching a Python Course to engineers and scientific
         people working on the ALMA project.}
        {Operations Support Facility (OSF), Atacama, Chile}
        {\myurl{http://www.almaobservatory.cl}}


%%%%%%%%%%%%%%%%%%%%%%%%%%%%%%%%%%%%%%%%%%%% Computer Science Skills
\section{\LARGE Computer Science skills}
\cvcomputer{\cc{Main Languages}}{C/C++, Python, Bash}
           {\cc{Secondary Languages}}{Fortran, Java, Octave/Matlab, R}
\cvcomputer{\cc{Parallel Computing}}{CUDA, OpenMPI, OpenMP, POSIX threads}
           {\cc{Libraries}}{SciPy, Pandas, Jupyter, ZeroMQ, cuBLAS, Intel MKL, Cython, Boost, SWIG}
%\cvcomputer{Operating Systems}{Linux Distributions and Windows}
\cvcomputer{\cc{SysAdmin}}{LAMP, Git, SVN, Jenkins}
           {\cc{WebDev}}{PHP, CSS, HTML, JavaScript, MySQL, Jekyll}
\cvcomputer{\cc{Scripting}}{Python, Bash, awk, sed}
           {\cc{Data bases}}{MySQL, PostgreSQL, MongoDB}
\cvcomputer{\cc{\textbf{Learning}}}{Scala, Go, AWS, Haskell}
           {}{}
           % Feature-driven development
           % Scrum
           % Kanban
%\cvcomputer{\cc{Methods}}{Agile development (FDD, Scrum, XP)}
%           {\cc{Practices}}{Continuous integration (CI)}

%%%%%%%%%%%%%%%%%%%%%%%%%%%%%%%%%%%%%%%%%%%% Languages
\section{\LARGE Languages}
\cvlanguage{\cc{Spanish}}{Native}{-}
\cvlanguage{\cc{English}}{Fluent}{Speaks, Reads, Writes, Technical and Colloquial}
\cvlanguage{\cc{German}}{Waystage or elementary (A2)}{Speaks, Reads, Writes}
\cvlanguage{\cc{Catalan}}{Competent (A1)}{Speaks, Reads, Writes}

\newpage
\section{\LARGE Events attendance}
\cventry{2016}
        {Astro-GR, What can we learn from the first GW detection? \emph{(Invited talk)}}
        {Centro de Ciencias de Benasque Pedro Pascual}
        {Benasque, Spain}
        {\myurl{http://astro-gr.org/astro-gr-benasque-2016-what-can-we-learn-from-first-gw-detection/}}
        {}

\cventry{2014}
        {IAU Symposium \#312 \emph{(Poster presentation)}}
        {GraviDy, a modular, GPU-based, direct-summation $N-$body integrator}
        {Star Clusters and Black Holes in Galaxies across Cosmic Time}
        {Beijing, China}
        {\myurl{http://silkroad.bao.ac.cn/iaus312}}

\cventry{2014}
        {Astro-GR, Gravitational waves and electromagnetic observations of dense
        stellar systems. \emph{(Invited talk)}}
        {INAF- Osservatorio Astronomico di Roma}
        {Frascati, Italy}
        {\myurl{http://members.aei.mpg.de/amaro-seoane/astro-gr-rome/}}
        {}

\cventry{2013}
        {3rd Aarseth Meeting, on stellar dynamics.}
        {Argelander-Institut für Astronomie (University of Bonn)}
        {Bonn, Germany}
        {}%\myurl{http://www.astro.uni-bonn.de/~sambaran/christmas_nbody2013/christmas_nbody2013.html}}
        {}

\cventry{2012}
        {Astro-GR, Black Hole Growth in the Universe \emph{(Invited talk)}}
        {GraviDy: A modular direct $N-$body GPU integrator}
        {National Astronomical Observatories (NAOC) of Chinese Academy of Sciences
        (CAS)}
        {Beijing, China}
        {\myurl{http://astro-gr.aei.mpg.de/Astro-GR@Beijing-2012}}

\cventry{2011}
        {First Argentine GPGPU Computing School for Scientific Applications}
        {Facultad de Matemáticas, Astronomía y Física}
        {Universidad Nacional de Córdoba, Argentina}
        {\myurl{http://www.famaf.unc.edu.ar/grupos/GPGPU/EscuelaGPGPU2011/}}
        {}

\cventry{2011}
        {High Performance Computing School}
        {(HPC School)}
        {CMCC, Universidad de la Frontera (UFRO)}
        {Temuco, Chile}
        {\myurl{http://www.cmcc.ufro.cl/hpc2011/}}

\cventry{2010}
        {Pan-American Advanced Studies Institute}
        {(PASI)}
        {Scientific Computing in the Americas: the challenge of massive parallelism}
        {UTFSM, Valparaíso, Chile}
        {\myurl{http://www.bu.edu/pasi/}}

\cventry{2009}
        {EVIC}
        {VI Latin-american summer school in Computational Intelligence, %
         III Latin-american Robotics summer meeting}
        {Universidad de Chile}
        {Santiago}
        {\myurl{http://www.uchile.cl}}

\cventry{2008, 2009}
        {ALMA Common Software Workshop}
        {Attendance to the general meeting and final course (Advanced) of the ALMA
        Common Software workshop, developed in joint with the Atacama Large
        Millimiter/Submillimeter Array project.}
        {UTFSM}
        {Valparaíso}
        {}

\cventry{2009}
        {Scientific Computing Advanced Training}
        {(SCAT)}
        {Final SCAT Meeting \& Satellite School}
        {UTFSM, Chile}
        {}

\cventry{2007-2009}
        {Encuentro Linux}
        {National gathering of Linux users}
        {Chile}
        {}
        {\myurl{http://encuentrolinux.cl}}

%\renewcommand{\listitemsymbol}{-}
\renewcommand*{\listitemsymbol}{\hspace{2em}\labelitemi~}
\section{\LARGE Honour and Awards}
\cvlistitem{Master studies scholarship in Chile (2012) given by the %
            ``Comisión Nacional de Investigación Científica y Tecnológica  %
            (CONICYT)'' to students with a good academic performance and
            participation in academic research projects (\myurl{www.conicyt.cl}).}
\vspace{0.2cm}
\cvlistitem{Initiation in Scientific Research Scholarship (PIIC) given by %
            ``Dirección General de Investigación y Postgrado UTFSM'',
            to students with academic research interest during Master
            studies. (\myurl{http://www.dgip.utfsm.cl/})}
\vspace{0.2cm}
\cvlistitem{Master studies scholarship, financial help,
            and 100\% fee payment exemption due to good academic performance.
            ``Dirección General de Investigación y Postgrado UTFSM''
            (\myurl{http://www.dgip.utfsm.cl/})}
\vspace{0.2cm}
\cvlistitem{Be a part of the \emph{Honour List}, to the best academic %
            performance among all the students in the {\utfsm}.
            Years 2007, 2008, 2009, 2010, 2011 and 2012.}

%\newpage
%\section{\LARGE References}
%    \cventry{}
%            {Pau Amaro-Seoane}
%            {Senior Scientist}
%            {\aei}
%            {\myurl{http://astro-gr.org}}
%            {\texttt{pau.amaro.seoane@gmail.com}}
%
%\vspace{0.3cm}
%    \cventry{}
%            {Horst Von Brand}
%            {Central Department of Computational Services Director}
%            {{\utfsm}}
%            {\texttt{vonbrand@inf.utfsm.cl}}
%            {}
%
%\vspace{0.3cm}
%    \cventry{}
%            {Cecilia Reyes}
%            {Informatics Department Adjoint Professor}
%            {{\utfsm}}
%            {\texttt{reyes@inf.utfsm.cl}}
%            {}

%\section{Publications}
% Publications from a BibTeX file
\nocite{*}
\bibliographystyle{unsrt}
\renewcommand\refname{\LARGE Publications}
\bibliography{latex/pub-en}
\end{document}

